\chapter{Bilan de mon parcours}

\section{Synthèse de 25 ans d'évolution professionnelle}

Mon parcours révèle une progression cohérente de l'expertise technique industrielle vers l'enseignement supérieur, structurée autour de trois piliers fondamentaux qui s'alignent avec les exigences du Master MEEF PIF.

\subsection{Les trois piliers de mon expertise}

\begin{enumerate}
\item \textbf{Expérience disciplinaire significative (20 ans d'industrie)}
   \begin{itemize}[leftmargin=*]
       \item Maîtrise étendue des technologies informatiques : développement, systèmes, réseaux, IA
       \item Veille technologique constante et adaptation aux évolutions
       \item Expérience de projets complexes (200K€ à 1,5M€)
   \end{itemize}

\item \textbf{Compétences pédagogiques développées (8 ans)}
   \begin{itemize}[leftmargin=*]
       \item Formation de 50+ collaborateurs en entreprise
       \item 5 ans d'enseignement NSI/SNT avec formation continue active
       \item Innovation pédagogique via la Forge des Communs Numériques
   \end{itemize}

\item \textbf{Démarche de recherche-action (3 ans)}
   \begin{itemize}[leftmargin=*]
       \item Expérimentation IA en éducation avec documentation méthodique
       \item Contribution à l'innovation pédagogique nationale
       \item Accompagnement de la communauté enseignante
   \end{itemize}
\end{enumerate}



\section{Valeur ajoutée de mon profil}

\textbf{Une triple compétence complémentaire :}
\begin{itemize}[leftmargin=*]
    \item \textbf{Technique :} 20 ans d'expertise industrielle couvrant tous les domaines du Master MEEF
    \item \textbf{Pédagogique :} Expérience managériale transférée vers l'enseignement avec succès
    \item \textbf{Recherche :} Démarche d'innovation et d'expérimentation déjà engagée
\end{itemize}

\textbf{Résultats observés :}
\begin{itemize}[leftmargin=*]
    \item 50+ professionnels formés en entreprise
    \item 10+ applications éducatives développées
    \item Rayonnement national via les initiatives ministérielles
\end{itemize}

\section{Cohérence et perspectives}

Ce parcours illustre la \textbf{complémentarité attendue} dans le Master MEEF entre expertise disciplinaire, compétences pédagogiques et démarche de recherche \parencite{referentiel_meef}. L'évolution naturelle vers l'enseignement supérieur s'appuie sur cette triple compétence acquise.

\textbf{Projet professionnel structuré :} Le Master MEEF constituera le tremplin académique nécessaire pour formaliser cette expertise et évoluer vers l'enseignement supérieur et la recherche en sciences de l'éducation \parencite{intelligence_artificielle_education,recherche_action_education}.

\section{Transition vers l'analyse des expériences}

Cette première partie a présenté qui je suis et d'où je viens. La partie suivante démontrera comment mes expériences correspondent concrètement aux compétences attendues dans le Master MEEF PIF (Parcours Formateur de Formateurs - FFMS).

L'analyse détaillée de mes expériences professionnelles révélera la \textbf{transférabilité directe} de mes compétences vers les missions d'un ingénieur de formation et enseignant-chercheur, en s'appuyant sur des exemples précis et des situations concrètes vécues sur le terrain.

\textbf{Objectif de la partie 2 :} Établir les correspondances explicites entre mes acquis professionnels et les compétences du référentiel MEEF PIF, en démontrant que mon parcours atypique constitue un atout majeur pour l'ingénierie de formation et la formation de formateurs.